\documentclass[10pt,a4paper]{amsart}
\usepackage[margin=19mm]{geometry}
\usepackage{amsmath,amssymb,mathtools}
\usepackage[scr=rsfs,cal=euler]{mathalfa}
\usepackage{xfrac}
\usepackage[unicode,hidelinks,bookmarks=false]{hyperref}
\newcommand{\ie}{i.e.\ }
\newcommand{\cf}{cf.\ }
\newcommand{\resp}{respectively\ }
\newcommand{\Subsection}{\subsection}
\providecommand{\nobreakdash}{\nobreak\hbox{-}}
\setlength{\emergencystretch}{2em}
\multlinegap0pt
\usepackage{bm}
\usepackage{bbm}
\usepackage{dsfont}
\usepackage{xspace}
\usepackage{cancel}
\usepackage{mathtools}
\usepackage{amsthm}
\makeatletter
\@ifundefined{definition}{\newtheorem{definition}{Definition}}{}
\@ifundefined{lemma}{\newtheorem{lemma}{Lemma}}{}
\@ifundefined{proposition}{\newtheorem{proposition}{Proposition}}{}
\@ifundefined{theorem}{\newtheorem{theorem}{Theorem}}{}
\makeatother
\newcommand{\R}{\mathbb{R}}
\title[On the asymptotic profile of solutions to semilinear damped ]{On the asymptotic profile of solutions to semilinear damped wave equations with critical nonlinearities}
\author{Trung Loc Tang, Dinh Van Duong}
\date{Source: arXiv:2509.19835v2 (CC BY 4.0)}
\begin{document}
\maketitle

\noindent\emph{Source and licence.} On the asymptotic profile of solutions to semilinear damped wave equations with critical nonlinearities. Version arXiv:2509.19835v2 (CC BY 4.0), distributed under the Creative Commons Attribution 4.0 International licence, as recorded in the arXiv licence metadata for this exact version and shown on the abstract page https://arxiv.org/abs/2509.19835v2. Full rights notice: licenses/arxiv-2509.19835-CC-BY-4.0.txt.\par\smallskip

\noindent\textit{Editorial scope.} Counted unit: the Proposition whose source label is \texttt{Pro1.1} in the article (source0.tex lines 417--423, source label \texttt{Pro1.1}), delivered together with the article's own complete original proof of it (source0.tex lines 424--476, one \texttt{proof} environment of 53 lines). No mathematics is written, repaired or re-derived: statement and proof are the article's own text line for line. Required context retained uncounted: Main.Eq.1 (lines 66--71); Condition1.1.1 (lines 131--140); Theorem1 (lines 165--189); LinearEstimates (lines 272--283); lemma1.3 (lines 380--386). Every definition, display and label of those lines is retained unchanged. Omitted from this package and not redistributed: 1--65, 72--130, 141--164, 190--271, 284--379, 387--416, 477--951. Nothing inside a retained excerpt is omitted or altered: every retained body is delivered exactly as the article's lines are written, so no omission receipt of any kind accompanies this package. Citations: the retained excerpts carry no citation command at all, so no bibliography entry of the article is transcribed and no reference is redistributed. Reference adaptation: none was needed and none was made; every reference inside the delivered unit resolves inside it, so no unresolved reference or citation is produced. Printed slips kept literal: no printed word, symbol, hypothesis or number of the counted statement or of its proof was altered. Layout and class: the article's own class is \texttt{amsart} with packages including amsmath, amsxtra, latexsym, amsthm, amssymb, amscd, pb-diagram, mathrsfs, mathabx, amsfonts, hyperref, geometry, bm, color; the wrapper is the lane's pinned \texttt{amsart} editorial wrapper at 10pt on A4, which restates the theorem-like environments the retained text needs and adds the article's own macro definitions that the retained text needs (\texttt{\textbackslash R}), with extra wrapper packages (bm, bbm, dsfont, xspace, cancel, mathtools). The delivered numbering is the unit's own. Assets: no retained span contains a figure environment, a graphics-inclusion command, a PDF-page inclusion, a table or a TikZ picture; no image file is copied into this package, the package declares no asset dependency and no image file is redistributed with it. Licence: version arXiv:2509.19835v2 of this article is distributed under the Creative Commons Attribution 4.0 International licence, as recorded in the arXiv licence metadata for this exact version and shown on the abstract page https://arxiv.org/abs/2509.19835v2. A full rights notice is delivered at licenses/arxiv-2509.19835-CC-BY-4.0.txt. Characters: every retained excerpt is pure ASCII, so the delivered engine, XeLaTeX with its default Latin Modern text font, reports no missing character.
\begin{equation} \label{Main.Eq.1}
\begin{cases}
u_{tt} -\Delta u + u_t= \mathcal{N}(u), &\quad x\in \R^n,\, t > 0, \\
(u, u_t)(0,x) = \varepsilon (u_0, u_1)(x), &\quad x\in \R^n, \\
\end{cases}
\end{equation}

\begin{definition}
    Let $\mu : [0, +\infty) \to [0, +\infty)$ be a modulus of continuity. Then, $\mu$ satisfies the Dini condition if
    \begin{align}
        \int_0^{1} \frac{\mu(s)}{s} ds < +\infty. \label{Condition1.1.1}
    \end{align}
    On the other hand, $\mu$ satisfies the non-Dini condition if (\ref{Condition1.1.1}) does not hold, i.e., if
    \begin{align}
        \int_0^{1} \frac{\mu(s)}{s} ds = +\infty. \label{Condition1.2.1}
    \end{align}
\end{definition}

\begin{theorem}[\textbf{Global existence}]\label{Theorem1}
     Let $1 \leq n \leq 4$. The modulus of continuity $\mu(s) $ satisfies the Dini condition (\ref{Condition1.1.1})
     and
     \begin{align}\label{Condition1.1.2}
         s|\mu'(s)| \lesssim |\mu(s)|  \quad \text{ for } s \in (0, 1].
     \end{align}
    In addition, we fix $$\alpha := \min\left\{2, 1+\frac{2}{n}\right\}, \,\,\beta_\alpha := (n-1)\left(\frac{1}{\alpha}-\frac{1}{2}\right)$$ and assume that
    \begin{align*}
        r \in \left(2, \frac{2n}{[n-2]^+}\right).
    \end{align*}
    Furthermore, the initial data $(u_0, u_1)$ satisfies
    \begin{align*}
        (u_0, u_1) \in \mathcal{D} := \left(H^2 \cap H^{2}_r \cap H^{\beta_\alpha}_\alpha\cap L^1  \right)\times \left(H^1 \cap H^1_r \cap L^{\alpha} \cap L^1\right).
    \end{align*}
    Then, there exists a constant $\bar{\varepsilon} > 0$ such that for any $\varepsilon \in (0, \bar{\varepsilon}]$, problem (\ref{Main.Eq.1}) admits a unique global (in time) Sobolev solution
    \begin{align*}
        u \in \mathcal{C}([0, \infty), H^2 \cap L^\alpha \cap L^\infty)
    \end{align*}
    satisfying the following estimates:
    \begin{align*}
        \|u(t,\cdot)\|_{\dot{H}^2} &\lesssim \varepsilon(1+t)^{-\frac{n}{4}-1} \|(u_0,u_1)\|_{\mathcal{D}},\\
        \|u(t,\cdot)\|_{L^{\alpha}} &\lesssim \varepsilon(1+t)^{-\frac{n}{2}(1-\frac{1}{\alpha})} \|(u_0,u_1)\|_{\mathcal{D}},\\
        \|u(t,\cdot)\|_{L^\infty} &\lesssim \varepsilon(1+t)^{-\frac{n}{2}}  \|(u_0,u_1)\|_{\mathcal{D}}.
    \end{align*}
\end{theorem}

\begin{lemma}[\textbf{Linear Estimates}]\label{LinearEstimates}
    Let $n \geq 1,\, j \in \{ 0,1\}, \,
    1 \leq \rho \leq q < \infty$, $q \ne 1$, $\beta_q := (n-1)\left|\frac{1}{2}-\frac{1}{q}\right|$ and $s_1 \geq s_2 \geq 0$.  Then, the following estimate holds for all $t > 0$:
    \begin{align*}
        &\|\partial_t^j |\nabla|^{s_1} \mathcal{K}(t,x)\ast_x \varphi(x)\|_{L^q}\\ &\quad\lesssim (1+t)^{-\frac{n}{2}(\frac{1}{\rho}-\frac{1}{q})-\frac{s_1-s_2}{2}-j} \||\nabla|^{s_2} \chi_{\rm L}(|\nabla|)\varphi\|_{L^{\rho}} + e^{-ct} \||\nabla|^{s_1}\chi_{\rm H}(|\nabla|) \varphi\|_{H_{q}^{\beta_q+j-1}},
        \end{align*}
        where $c$ is a suitable positive constant.
    Furthermore, we obtain the following estimate for $ 1\leq \rho \leq \infty, 1 < r < \infty$ and $d > \displaystyle\frac{n}{r}$:
    \begin{align*}
        &\|\partial_t^j |\nabla|^{s_1}\mathcal{K}(t,x) \ast_x \varphi(x) \|_{L^{\infty}}\\ &\quad\lesssim (1+t)^{-\frac{n}{2\rho}-\frac{s_1-s_2}{2}-j} \||\nabla|^{s_2} \chi_{\rm L}(|\nabla|) \varphi\|_{L^\rho} + e^{-ct} \||\nabla|^{s_1 + d} \chi_{\rm H}(|\nabla|) \varphi\|_{H^{\beta_r+j -1}_r}.
    \end{align*}
\end{lemma}

\begin{lemma}\label{lemma1.3}
    Under the assumptions of Theorem \ref{Theorem1}, the following estimates hold for all $u \in X(T,\varepsilon_0)$ and $T >0, \,\, \varepsilon_0 \in (0, 1]$:
    \begin{align*}
    \|\mathcal{N}(u(\tau,\cdot))\|_{L^{\gamma}} &\lesssim \mu\left(\varepsilon_0 (1+\tau)^{-\frac{n}{2}}\right) (1+\tau)^{-\frac{n}{2}(1+\frac{2}{n}-\frac{1}{\gamma})}  \|u\|_{X(T)}^{1+\frac{2}{n}} \quad\text{ for } \gamma \geq 1,\\
    \|\nabla \mathcal{N}(u(\tau,\cdot))\|_{L^{\theta}} &\lesssim \mu\left(\varepsilon_0 (1+\tau)^{-\frac{n}{2}}\right) (1+\tau)^{-\frac{n}{2}(1+\frac{2}{n}-\frac{1}{\theta})-\frac{1}{2}}  \|u\|_{X(T)}^{1+\frac{2}{n}} \quad \text{ for } \theta \in \left[2, \frac{2n}{[n-2]^+}\right).
    \end{align*}
\end{lemma}

\begin{proposition}\label{Pro1.1}
     Under the assumptions of Theorem \ref{Theorem1}, the following estimate holds for all $u \in X(T, \varepsilon_0),\, \varepsilon_0 \in (0, 1]$:
     \begin{align*}
         \|u^{\rm non}\|_{X(T)} \lesssim \mathcal{I}(\varepsilon_0)\|u\|_{X(T)}^{1+\frac{2}{n}}.
     \end{align*}
   
\end{proposition}

     \begin{proof}
         Firstly, using Lemma \ref{lemma1.3} and
         Lemma \ref{LinearEstimates} with $ j= 0, \,\rho = 1, q = 2$, $(s_1, s_2) =(2,0)$ for $\tau \in [0, t/2]$, and $\rho=q =2, (s_1,s_2) =(2,1)$ for $\tau \in (t/2, t]$ to derive
         \begin{align*}
             \| u^{\rm non}(t,\cdot)\|_{\dot{H}^2} &\lesssim \int_0^{t/2} (1+t-\tau)^{-\frac{n}{4}-1} \| \mathcal{N}(u(\tau,\cdot))\|_{L^1 \cap \dot{H}^1} d\tau\\
             &\qquad + \int_{t/2}^t (1+t-\tau)^{-\frac{1}{2}} \| \mathcal{N}(u(\tau,\cdot))\|_{ \dot{H}^1} d\tau\\
             &\lesssim \|u\|_{X(T)}^{1+\frac{2}{n}}\int_0^{t/2} (1+t-\tau)^{-\frac{n}{4}-1} (1+\tau)^{-1} \mu\left(\varepsilon_0 (1+\tau)^{-\frac{n}{2}}\right) d\tau\\
             &\qquad + \|u\|_{X(T)}^{1+\frac{2}{n}}\int_{t/2}^t (1+t-\tau)^{-\frac{1}{2}} (1+\tau)^{-\frac{n}{4}-\frac{3}{2}} \mu\left(\varepsilon_0 (1+\tau)^{-\frac{n}{2}}\right) d\tau\\
             &\lesssim  (1+t)^{-\frac{n}{4}-1}\|u\|_{X(T)}^{1+\frac{2}{n}} \int_0^{t/2}(1+\tau)^{-1} \mu\left(\varepsilon_0 (1+\tau)^{-\frac{n}{2}}\right) d\tau\\
             &\qquad +  (1+t)^{-\frac{n}{4}-1}\|u\|_{X(T)}^{1+\frac{2}{n}} \int_{t/2}^t (1+t-\tau)^{-1} \mu\left(\varepsilon_0 (1+t-\tau)^{-\frac{n}{2}}\right) d\tau,
         \end{align*}
where noting that $1+t \sim 1+\tau \geq 1+t-\tau$ for all $\tau \in [t/2, t]$. Using the change of variables $\theta_1 = \varepsilon_0 (1+\tau)^{-\frac{n}{2}}$ and $\theta_2 = \varepsilon_0(1+t-\tau)^{-\frac{n}{2}}$, we get
\begin{align}
    \int_0^{t/2} (1+\tau)^{-1} \mu\left(\varepsilon_0 (1+\tau)^{-\frac{n}{2}}\right) d\tau   +
    \int_{t/2}^{t} (1+t-\tau)^{-1} \mu\left(\varepsilon_0 (1+t-\tau)^{-\frac{n}{2}}\right) d\tau \lesssim \mathcal{I}(\varepsilon_0). \label{Impor_Re1}
\end{align}
Therefore, we can conclude that 
    \begin{align}
        \|u^{\rm non}(t,\cdot)\|_{\dot{H}^2} \lesssim \mathcal{I}(\varepsilon_0)(1+t)^{-\frac{n}{4}-1} \|u\|_{X(T)}^{1+\frac{2}{n}}. \label{Main.Es.1}
    \end{align}
 Next, thanks to again 
         Lemma \ref{LinearEstimates} for $j =0, \,\rho = 1, q = \infty$, $(s_1, s_2) = (0,0)$ for $\tau \in [0, t/2]$, and $\rho =r, q= \infty, \,(s_1, s_2) =(0,0)$ for $\tau \in (t/2, t]$ to obtain
    \begin{align*}
        \|u^{\rm non}(t,\cdot)\|_{L^\infty} &\lesssim \int_0^{t/2} (1+t-\tau)^{-\frac{n}{2}} \|\mathcal{N}(u(\tau,\cdot))\|_{L^1 \cap H_r^{d+\beta_r-1}} d\tau\\
        &\qquad +\int_{t/2}^t (1+t-\tau)^{-\frac{n}{2r}} \|\mathcal{N}(u(\tau,\cdot))\|_{L^r \cap H_r^{d+\beta_r-1}} d\tau,
    \end{align*}
    where $d := n/r +\delta$ with $\delta$ is a sufficiently small constant. 
Noting that the conditions $r \in \left(2, \displaystyle\frac{2n}{[n-2]^+}\right)$ and $1 \leq n \leq 4$ imply
\begin{align}
    0 < d + \beta_r = \frac{n}{r} +\delta + (n-1)\left(\frac{1}{2}-\frac{1}{r}\right) < 2,  \label{Re10}
\end{align}
   that is, 
   \begin{align*}
       \| \mathcal{N}(u(\tau,\cdot))\|_{H_r^{d+\beta_r-1}} &\lesssim \| \mathcal{N}(u(\tau,\cdot))\|_{H_r^1} \sim \|\mathcal{N}(u(\tau,\cdot))\|_{L^r} + \|\nabla \mathcal{N}(u(\tau,\cdot))\|_{L^r}\\ 
       &\lesssim (1+\tau)^{-\frac{n}{2}(1+\frac{2}{n}-\frac{1}{r})} \mu\left(\varepsilon_0 (1+\tau)^{-\frac{n}{2}}\right) \|u\|_{X(T)}^{1+\frac{2}{n}},
   \end{align*}
   due to using again Lemma \ref{lemma1.3}. Combining this with the relation (\ref{Impor_Re1}) and $n < 2r$, we get
   \begin{align}
       \|u^{\rm non}(t,\cdot)\|_{L^\infty} &\lesssim \|u\|_{X(T)}^{1+\frac{2}{n}}\int_0^{t/2} (1+t-\tau)^{-\frac{n}{2}} (1+\tau)^{-1} \mu\left(\varepsilon_0(1+\tau)^{-\frac{n}{2}}\right) d\tau \notag\\
       &\quad+ \|u\|_{X(T)}^{1+\frac{2}{n}} \int_{t/2}^t (1+t-\tau)^{-\frac{n}{2r}} (1+\tau)^{-\frac{n}{2}(1+\frac{2}{n}-\frac{1}{r})} \mu\left(\varepsilon_0 (1+\tau)^{-\frac{n}{2}}\right) d\tau \notag\\
       &\lesssim (1+t)^{-\frac{n}{2}} \|u\|_{X(T)}^{1+\frac{2}{n}} \int_0^{t/2} (1+\tau)^{-1} \mu\left(\varepsilon_0 (1+\tau)^{-\frac{n}{2}}\right) d\tau \notag\\
       &\quad + (1+t)^{-\frac{n}{2}} \|u\|_{X(T)}^{1+\frac{2}{n}} \int_{t/2}^t (1+t-\tau)^{-\frac{n}{2r}} (1+\tau)^{-1+\frac{n}{2r}}  \mu\left(\varepsilon_0 (1+\tau)^{-\frac{n}{2}}\right) d\tau \notag\\
       &\lesssim \mathcal{I}(\varepsilon_0) (1+t)^{-\frac{n}{2}} \|u\|_{X(T)}^{1+\frac{2}{n}}. \label{Main.Es.2}
   \end{align}
   Finally, we need to estimate $\|u^{\rm non}(\tau,\cdot)\|_{L^{\alpha}}$. We have $\beta_\alpha < 1$ for all $1 \leq n \leq 4$. Therefore, employing again Lemma \ref{lemma1.3} combined with Lemma \ref{LinearEstimates} with $j = 0, \,\rho = 1, q =\alpha, (s_1, s_2) = (0,0)$ for $\tau \in [0, t/2)$ and $\rho = q =\alpha, (s_1, s_2) = (0,0)$ for $\tau \in [t/2, t]$ we have
   \begin{align}
       \|u^{\rm non}(t,\cdot)\|_{L^\alpha} &\lesssim \int_0^{t/2} (1+t-\tau)^{-\frac{n}{2}(1-\frac{1}{\alpha})} \|\mathcal{N}(u(\tau,\cdot))\|_{L^1 \cap L^\alpha} d\tau  +\int_{t/2}^t \|\mathcal{N}(u(\tau,\cdot))\|_{L^{\alpha}} d\tau \notag\\
       &\lesssim \|u\|_{X(T)}^{1+\frac{2}{n}} \int_0^{t/2} (1+t-\tau)^{-\frac{n}{2}(1-\frac{1}{\alpha})} (1+\tau)^{-1} \mu\left(\varepsilon_0 (1+\tau)^{-\frac{n}{2}}\right) d\tau \notag\\
       &\quad+ \|u\|_{X(T)}^{1+\frac{2}{n}}\int_{t/2}^t (1+\tau)^{-\frac{n}{2}(1+\frac{2}{n}-\frac{1}{\alpha})} \mu\left(\varepsilon_0 (1+\tau)^{-\frac{n}{2}}\right) d\tau \notag\\
       &\lesssim \mathcal{I}(\varepsilon_0) (1+t)^{-\frac{n}{2}(1-\frac{1}{\alpha})} \|u\|_{X(T)}^{1+\frac{2}{n}}. \label{Main.Es.3}
   \end{align}
   The estimates (\ref{Main.Es.1})-(\ref{Main.Es.3}) complete the proof of Proposition \ref{Pro1.1}.
     \end{proof}

\end{document}
